\chapter{LLM Agent}
\label{ch:llm}

This second agent is a cognitive coordinator. It is built on the same BDI core as the
executor, owning a \texttt{BeliefBase} and an \texttt{IntentionEngine}, so it can
itself move, pick up, and deliver, including the PDDL-based crate clearing described in 
\ref{ch:pddl}. 
What sets it apart is an additional natural-language reasoning layer. Its job is to 
receive instructions from the Admin in the chat, reason
about them, and translate them into concrete commands and policies for the whole
team. It never invents game actions from this reasoning directly: every effect is
mediated either by a tool or by a peer-to-peer message, and any movement it commands,
whether for itself or the partner, is ultimately carried out by a BDI agent through
the loop of the previous chapters.
There were 3 possible models to chose from and the one selected is \texttt{llama-3.3-70b-lmstudio}.

\begin{center}
\begin{verbatim}
  Admin prompt --> reasoning loop --> tool call --> {local effect, P2P message}
                        ^                  |
                        +---- tool result -+        (repeat until "stop")
\end{verbatim}
\end{center}

\section{The reasoning loop}

A prompt triggers an iterative, single-step reasoning loop driven by a language
model under a fixed system prompt. The model returns one validated JSON
object per turn, of exactly one of three kinds:

\begin{itemize}
  \item \textbf{tool}, a request to execute one of the available tools;
  \item \textbf{answer}, a textual output to return;
  \item \textbf{stop}, a signal that the prompt has been fully handled.
\end{itemize}

Chain-of-thought is elicited through delimited reasoning that is stripped before
parsing, and the model runs at zero temperature for determinism. After each turn the
model executes the requested tool (or records the answer), feeds the result
back as the next input, and repeats until a \texttt{stop} is produced. Outputs that
violate the schema are rejected and re-requested with an error description, up to three attempts. 
The loop advances one step at a time by design, so a message containing
several tasks \uline{is decomposed and resolved in presentation order}, with intermediate
answers returned in that same order(\ref{sub:res_and_lim}).




\section{Prompt construction}

The system prompt is fixed for the whole run and organised into clearly delimited
sections. A \emph{role} description frames the model as the team's reasoning brain
and lists the agent identifiers. A \emph{definitions} block gives precise meanings
for the three concepts the coordinator must tell apart: a \emph{reward} (the points
gained or lost by carrying out an instruction, possibly written as an expression), a
\emph{task} (a single unit of work, where one message may contain several), and a
\emph{rule} (a standing instruction that changes how future tasks are scored). A set
of operating \emph{rules} then fixes the procedure the model must follow, and a
\emph{response-format} contract pins down the exact JSON it may emit. The prompt ends
with the available tools and a bank of worked examples.

The definitions and operating rules are what make the behaviour reliable rather than
improvised. Because \emph{reward}, \emph{task}, and \emph{rule} are defined
explicitly, the model can consistently separate a one-off rewarded task from a
standing rule, the distinction that drives feasibility gating (a task is checked for
a positive reward before acting, a rule is simply applied). The operating rules
encode the rest of the discipline: expressions must be solved with the arithmetic
tool instead of computed by the model, multi-task prompts must be decomposed and
answered in presentation order, and unknown facts such as positions or map extremes
must be fetched with a context tool rather than guessed. They also carry targeted
guardrails learned from failures, for instance checking that a target tile is
walkable before a move is ordered, and never sending both agents to the same tile.

Beyond stating the rules, the prompt teaches them largely by example. It is
few-shot, with curated cases covering feasibility checks, multi-task decomposition,
rule application, holds and resumes, and cooperation contracts, each showing the full
turn-by-turn exchange. \texttt{The examples, more than the rules alone, are what holps the model to 
hold the JSOM protocol and execute one step at a time.(remove)}
Finally, the tools section of the prompt
is generated automatically from the same registry the coordinator executes against,
so the descriptions the model reads can never drift from the tools that actually run.


\section{Message filtering and history}

Two filters precede the reasoning loop, both enforced in code rather than left to the
prompt. First, only messages originating from the Admin are treated as candidate
instructions; messages from any other agent are routed to the P2P handler instead.
Second, the coordinator distinguishes actionable instructions from Q\&A messages
carrying no reward: a rewardless, non-command message produces a \texttt{stop} with
no chat output at all. Conversation history is kept in full but passed to the model
only when needed and in filtered form, namely past questions and answers only, never
the intermediate reasoning or tool calls, so the prompt stays compact.

\section{Feasibility gating}

The central reasoning discipline is that effectful actions are gated on a feasibility
check. Any expression, whether a reward value or a logical condition, must be resolved
through the arithmetic tool rather than computed by the model itself. World- or
rule-changing actions (movement, variable setting, cooperation proposals) are
reward-gated: they are issued only once the arithmetic evaluation has confirmed a
strictly positive reward for the current prompt. Once confirmed, the gate is sticky
for the rest of that prompt, so later numeric evaluations such as computing
coordinates do not re-lock it. A prompt with no reward yields no action and no chat
output.

Two classes are intentionally exempt. \emph{Standing policy rules}
(``from now on\ldots'', penalties, multipliers) are announcements of an effect, not
rewarded tasks, so they are always applied regardless of reward, since they describe
how the team adapts to penalties. \emph{Control and query} actions (hold, resume,
context and history lookup, contract cancellation) are likewise exempt, as they carry
no reward of their own.

\section{Tools}

Tools fall into two groups. \emph{Read} tools retrieve facts without side effects:
arithmetic evaluation, the local context snapshot (own and peer state, parcels, map
extremes), and conversation history. \emph{Act} tools change behaviour: directing an
agent to a coordinate, applying scoring or avoidance policy rules, setting shared
variables, holding or resuming agents, and proposing cooperation contracts.

The key design point is symmetry of effect. When an act tool targets the coordinator
itself (or ``all''), it mutates the coordinator's own beliefs directly; when it
targets the executor (or ``all''), it additionally emits a peer-to-peer message so
the executor applies the same change to its belief base. The same vocabulary governs
both agents through one uniform mechanism, detailed in Chapter \ref{ch:coordination}.
